\section{Conclusion}
\label{sec:conclusion}

We decompose transformer updates into parallel and perpendicular components.
Across pretrained models, targeted edits show that robustness depends on both
the representation space and reference direction: self-value-aligned attention
components are comparatively stable, whereas residual-space and perpendicular
edits are more fragile. Compression diagnosis and value-parallel pretraining
demonstrate supporting uses of this geometry.

\section*{Limitations}
\label{sec:limitations}

Our evidence is limited to the evaluated decoder-only models, text tasks,
intervention sites, and training budgets; the same sensitivities may not hold
for other architectures, modalities, long-context distributions, or fully
trained models. The decomposition is also reference- and site-dependent:
attention supplies a natural current-token value reference, whereas an MLP has
no canonical internal analogue, making the appendix study a scoped diagnostic.
Finally, these edits are analysis tools rather than deployment-ready methods.
The matched-seed training study improves early validation loss without a stable
downstream gain, so we make no general efficiency or performance claim.

\section*{Ethical Considerations}

Directional decomposition provides a concrete tool for auditing representation
change and screening compression or interventions before downstream evaluation.
This screening requires no retraining. Average stability does not imply safety:
components can remain essential for particular prompts, tasks, or contexts, and
small edits can silently alter behavior. Geometry cannot replace task-specific
testing or human review. We do not evaluate bias, privacy, security, or
high-stakes use, and these tools could also enable targeted manipulation.
